\begin{note}
    $x, y \in C_G(A)$ does not necessarily mean $xy = yx$.
\end{note}
\begin{definition}[Center]
    \label{def:center}
    Define $Z(G) = \{g \in G: gx = xg \text{ for all } x \in G\}$, the set of elements commuting with all the elements of $G$. This subset of $G$ is called the $center$ of $G$.     
\end{definition}
\begin{note}
    $Z(G) = C_G(G)$, so the argument above proves $Z(G) \le G$ as a special case.   
\end{note}
\begin{definition}[Normalizer]
    Let $g A g^{-1} := \{g a g^{-1} : a \in A\}$. The \textit{normalizer} of $A \in G$ is the set $N_g(A) = \{g  \in G : g A g^{-1} = A\}$.  
\end{definition}
$N_G(A)$ is a subgroup of $G$ (follows from the same steps which demonstrated that $C_G(A) \le G$ with appropriate modifications).
\section{Cayley's Table}
\textbf{To be written...}